% Part: incompleteness
% Chapter: representability-in-q
% Section: representable-comp

\documentclass[../../../include/open-logic-section]{subfiles}

\begin{document}

\olfileid{inc}{req}{rpc}
\olsection{Fungsi Representabel ing $\Th{Q}$ Iku Komputabel}

Kita bakal mbuktekake yen saben fungsi kang representabel ing~$\Th{Q}$
iku komputabel. Luwih dhisik, kita kudu netepake sawijining lema
ngenani fungsi kang representabel ing~$\Th{Q}$.

\begin{lem}\ollabel{lem:rep-q}
  Yen $f(x_0,
  \dots, x_k)$ representabel ing~$\Th{Q}$, ana
  !!a{formula}~$!A_f(x_0, \dots, x_k, y)$ saengga
  \[
  \Th{Q} \Proves !A_f(\num{n_0}, \dots, \num{n_k}, \num{m})
  \quad\text{yen lan mung yen}\quad m = f(n_0, \dots, n_k).
  \]
\end{lem}

\begin{proof}
  Arah ``yen'' iku
\olref[int]{defn:representable-fn}\olref[int]{defn:rep:a}. Arah ``mung
yen'' katon kaya mangkene: Upamana $\Th{Q} \Proves !A_f(\num{n_0},
\dots, \num{n_k}, \num{m})$ nanging $m \neq f(n_0, \dots, n_k)$. Tetepna $l =
f(n_0, \dots, n_k)$. Miturut
\olref[int]{defn:representable-fn}\olref[int]{defn:rep:a}, $\Th{Q}
\Proves !A_f(\num{n_0}, \dots, \num{n_k}, \num{l})$. Miturut
\olref[int]{defn:representable-fn}\olref[int]{defn:rep:b},
$\lforall[y][(!A_f(\num{n_0}, \dots, \num{n_k}, y) \lif \num{l} =
y)]$. Kanthi logika lan asumsi yen $\Th{Q} \Proves
!A_f(\num{n_0}, \dots, \num{n_k}, \num{m})$, kita entuk $\Th{Q}
\Proves \eq[\num{l}][\num{m}]$. Nanging, miturut
\olref[bre]{lem:q-proves-neq}, $\Th{Q} \Proves
\eq/[\num{l}][\num{m}]$. Dadi $\Th{Q}$ ora konsisten. Iki
ora bisa kelakon, amarga model baku nyukupi $\Th{Q}$ (delengen
\olref[int][def]{def:standard-model}), $\Sat{N}{\Th{Q}}$, lan
teori kang bisa dicukupi mesthi konsisten miturut Teorema Kesahihan
(\tagrefs{prfAX/{fol:axd:sou:cor:consistency-soundness},
prfSC/{fol:seq:sou:cor:consistency-soundness},
prfND/{fol:ntd:sou:cor:consistency-soundness},
prfTab/{fol:tab:sou:cor:consistency-soundness}}).
\end{proof}

\begin{lem}
Saben fungsi kang representabel ing $\Th{Q}$ iku komputabel.
\end{lem}

\begin{proof}
Luwih dhisik, ayo deleng gagasan intuitif ngapa pratelan iki bener.
Kanggo ngitung~$f$, kita nindakake iki. Dhaptarna kabeh
!!{derivation}s~$\delta$ kang bisa ana ing basa aritmetika. Iki
bisa ditindakake kanthi mekanis. Kanggo saben derivasi, priksa apa iku
!!a{derivation} saka !!a{formula} kang awujud~$!A_f(\num{n_0}, \dots,
\num{n_k}, \num{m})$ (!!{formula} kang makili $f$ ing~$\Th{Q}$
saka \olref{lem:rep-q}). Yen iya, $m = f(n_0, \dots, n_k)$ miturut
\olref{lem:rep-q}, lan kita wis nemokake bijine~$f$. Panggolekan iki
mandheg amarga $\Th{Q} \Proves !A_f(\num{n_0}, \dots, \num{n_k},
\num{f(n_0, \dots, n_k)})$, mula pungkasane kita nemokake $\delta$ kang
cocog.

Iki durung persis, amarga prosedur kita nggarap
!!{derivation}s lan !!{formula}s, dudu mung wilangan, lan kita
durung nerangake kanthi tepat ngapa ``ndhaptar kabeh !!{derivation}s
kang bisa ana'' bisa ditindakake kanthi mekanis. Nanging, kaya kang wis
kita deleng, suku, !!{formula}s, lan !!{derivation}s bisa dikodeake
nganggo wilangan G\"odel. Kita uga wis ngenalake model komputasi kang
persis, yaiku fungsi rekursif umum. Lan kita wis weruh yen relasi
$\Prf[\Th{Q}](d,y)$, kang lumaku yen lan mung yen $d$ iku wilangan
G\"odel saka !!a{derivation} kanggo !!{formula} kanthi wilangan
G\"odel~$y$ saka aksioma~$\Th{Q}$, iku rekursif (primitif). Fungsi
rekursif primitif liyane kang dibutuhake yaiku $\fn{num}$
(\olref[art][trm]{prop:num-primrec}) lan $\fn{Subst}$
(\olref[art][sub]{prop:subst-primrec}). Saka fungsi lan relasi iki,
$f$ bisa ditegesi liwat panggolekan tanpa wates; mula $f$ rekursif. Cathetan OLPL-187 lan SOL6-F082: kanggo kalkulus aksiomatik, Prf saiki nguji sertifikat kondisional tanpa asumsi kanthi anteseden winates saka Q, miturut pratelan kang wis dibenerake. Anane sertifikat kanggo target padha karo kabukten saka Q, mula panggolekan iki tetep menehi nilai fungsi. Kode sertifikat ora kudu padha kode derivasi asli saka Q.

Luwih dhisik, tegesna
\begin{multline*}
  A(n_0, \dots, n_k, m) = \\
  \fn{Subst}(\fn{Subst}(\dots\fn{Subst}(\Gn{!A_f}, \fn{num}(n_0), \Gn{x_0}),\\ \dots),
  \fn{num}(n_k),  \Gn{x_k}), \fn{num}(m), \Gn{y})
\end{multline*}
Iki katon ruwet, nanging mung fungsi $A(n_0, \dots, n_k,
m) = \Gn{!A_f(\num{n_0}, \dots, \num{n_k}, \num{m})}$.

Saiki, delengen relasi~$R(n_0, \dots, n_k, s)$ kang lumaku yen
$(s)_0$ iku wilangan G\"odel saka !!a{derivation} saka~$\Th{Q}$ kanggo
$!A_f(\num{n_0}, \dots, \num{n_k}, \num{(s)_1})$:
\[
R(n_0, \dots, n_k, s) \quad\text{yen lan mung yen}\quad \Prf[\Th{Q}]((s)_0, A(n_0,
\dots, n_k, (s)_1))
\]
Yen kita bisa nemokake~$s$ saengga $R(n_0, \dots, n_k, s)$ lumaku, kita wis
nemokake pasangan wilangan---$(s)_0$ lan~$(s)_1$---saengga $(s)_0$ iku
wilangan G\"odel saka !!a{derivation} kanggo~$A_f(\num{n_0}, \dots,
\num{n_k}, \num{(s)_1})$. Mula nggoleki~$s$ padha karo nggoleki pasangan
$d$ lan $m$ ing bukti informal mau. Fungsi komputabel kang
``nggoleki'' $s$ kaya mangkono bisa ditegesi liwat panggolekan tanpa
wates reguler. Gatekna yen $R$ reguler: kanggo saben $n_0$, \dots, $n_k$, ana
!!a{derivation}~$\delta$ kanggo $\Th{Q} \Proves !A_f(\num{n_0}, \dots,
\num{n_k}, \num{f(n_0, \dots, n_k)})$, mula $R(n_0, \dots, n_k, s)$
lumaku kanggo $s = \tuple{\Gn{\delta}, f(n_0, \dots, n_k)}$. Dadi, kita bisa
nulis $f$ dadi
\[
f(n_0,\dots,n_{k}) = (\umin{s}{R(n_0, \dots, n_k, s)})_1.
\]
\end{proof}

\end{document}
